\section{Conclusion}
\label{sec:conclusion}

We began by observing a coverage paradox: pylint flags 100\% of HumanEval baseline failures, yet pylint-guided repair \emph{reduces} HumanEval pass@1 by 4.3 percentage points. Our study shows that this paradox resolves cleanly once coverage is decomposed by flag category. The 100\% coverage is a mirage---it is driven by Convention-category style rules (C0304: missing newline, C0114: missing docstring) that fire universally on LLM-generated code regardless of functional correctness. Only 12.5\% of HumanEval failures receive a functional (Error or Warning category) flag; mypy coverage is 0\%. When the LLM receives ``fix your formatting'' as its primary repair guidance on a failed algorithmic solution, it may restructure code while addressing style---and on complex HumanEval problems, this rewriting introduces new logical errors.

Our main contributions are: (1) first iso-compute comparison confirming execution feedback significantly dominates pylint/mypy on HumanEval and MBPP (McNemar $p$=0.0001 and $p$<10$^{-18}$); (2) style-function dissociation measurement (94.3\% C-category, 12.5\% E+W functional coverage); (3) benchmark asymmetry finding revealing task-complexity moderation of static analysis feedback utility.

\textbf{Future directions.} Qwen2.5-Coder-7B replication; per-problem solution analysis for regression cases; W+E-only pylint filtering (removing C-category noise); larger budget comparison ($B$=2000+) for multi-round characterization; extension to code-specialized models and non-Python benchmarks.

For LLM code repair, the question is not whether a tool detects errors---it is whether it detects the \emph{right} errors. Pylint detects formatting; execution detects failure. At $B$=1000 on functional correctness benchmarks, the difference is +40.2pp on MBPP.
